\documentclass[11pt]{article}

\usepackage{amsmath,amssymb,bm}
\usepackage{hyperref}
\usepackage{svg}
\usepackage[a4paper,margin=1in]{geometry}
\usepackage{array}

\newcolumntype{C}{>{\centering\arraybackslash}m{0.21\textwidth}}

\title{Single-Target FIM Comparison:\\
OFDM vs OTFS, with and without OFDM-side ICI}
\author{}
\date{}

\begin{document}
\maketitle

\section{Setup}

We compare the single-target Fisher information matrices derived in
ofdm\_fim.tex (Section~3.1) and otfs\_fim.tex (Section~3.1) for the
four combinations
\begin{enumerate}
    \item OFDM with no ICI (subcarriers treated as orthogonal,
          \(I(n'-n) = \delta[n'-n]\)),
    \item OFDM with the full Doppler-induced ICI kernel
          \(I(n'-n) = (e^{j2\pi\nu T}-1)/(j2\pi[(n'-n)\Delta f + \nu])\),
    \item Two-step OTFS using the no-ICI form of the inner OFDM
          modulator,
    \item Two-step OTFS using the full ICI form of the inner OFDM
          modulator.
\end{enumerate}
In all four cases the parameter vector is
\(\bm{\eta} = [\tilde{\alpha}_R,\tilde{\alpha}_I,\tau,\nu]^T\), the
noise per observation bin has variance \(\sigma^2\), and the data
symbols are i.i.d.\ zero-mean unit-energy QPSK,
\begin{equation}
    \mathrm{E}\{X[n,m]\,X^*[n',m']\} = \delta_{n,n'}\,\delta_{m,m'},
    \qquad
    \mathrm{E}\{x[l,k]\,x^*[l',k']\} = \delta_{l,l'}\,\delta_{k,k'}.
\end{equation}
For OFDM with no ICI we have \(|X[n,m]|^2 = 1\) deterministically (no
expectation needed); in the other three cases cross-bin products
survive in the FIM integrand and only vanish on average, so we report
the expected FIM \(\mathrm{E}\{\bm{J}(\bm{\eta})\}\).

\subsection*{Notation}

Let \(P_a := \sum_{n=0}^{N-1} n^a\) and \(Q_b := \sum_{m=0}^{M-1} m^b\)
denote the elementary power sums on the OFDM grid, with
\(P_0 = N,\, P_1 = N(N-1)/2,\, P_2 = N(N-1)(2N-1)/6\) and analogously
for \(Q_b\). The OFDM-side ICI introduces the kernel-summed moments
\begin{align}
    \mathcal{S}_a &:= \sum_{n,n'=0}^{N-1} (n')^a\, |I(n'-n)|^2,
        \qquad a = 0,1,2, \\
    \mathcal{T}_a &:= \sum_{n,n'=0}^{N-1} (n')^a\, I^*(n'-n)\, I'(n'-n),
        \qquad a = 0,1, \\
    \mathcal{U}   &:= \sum_{n,n'=0}^{N-1} |I'(n'-n)|^2,
\end{align}
where \(I'(n) = T e^{j2\pi\nu T}/(j2\pi(n\Delta f + \nu))
                 - (e^{j2\pi\nu T} - 1)/(j2\pi(n\Delta f + \nu))^2\)
is the \(\nu\)-derivative kernel of ofdm\_fim.tex.
\(\mathcal{S}_a\) and \(\mathcal{U}\) are real; \(\mathcal{T}_a\) is
generally complex. In the limit \(I\to\delta\) we have
\(\mathcal{S}_a\to P_a\), \(\mathcal{T}_a\to 0\), \(\mathcal{U}\to 0\).

\section{Side-by-side comparison}

The expected FIM entries for all four cases are collected in
Table~\ref{tab:fim_compare}. Off-diagonal entries follow from symmetry
\([\bm{J}]_{k,l} = [\bm{J}]_{l,k}\); the \(4\times 4\) FIM has 10
distinct entries.

\begin{table}[h]
\centering
\resizebox{\textwidth}{!}{$
\renewcommand{\arraystretch}{2.4}
\begin{array}{l|c|c|c|c}
\hline
\text{FIM entry} & \text{OFDM, no ICI} & \text{OFDM, ICI} &
\text{OTFS, no OFDM-ICI} & \text{OTFS, OFDM-ICI} \\
\hline
J_{\tilde{\alpha}_R, \tilde{\alpha}_R} = J_{\tilde{\alpha}_I, \tilde{\alpha}_I}
    & \dfrac{2 N M}{\sigma^2}
    & \dfrac{2 M\, \mathcal{S}_0}{\sigma^2}
    & \dfrac{2 N M}{\sigma^2}
    & \dfrac{2 M\, \mathcal{S}_0}{\sigma^2} \\
J_{\tilde{\alpha}_R, \tilde{\alpha}_I}
    & 0 & 0 & 0 & 0 \\
J_{\tau, \tau}
    & \dfrac{8\pi^2 |\tilde{\alpha}|^2 M (\Delta f)^2\, P_2}{\sigma^2}
    & \dfrac{8\pi^2 |\tilde{\alpha}|^2 M (\Delta f)^2\, \mathcal{S}_2}{\sigma^2}
    & \dfrac{8\pi^2 |\tilde{\alpha}|^2 M (\Delta f)^2\, P_2}{\sigma^2}
    & \dfrac{8\pi^2 |\tilde{\alpha}|^2 M (\Delta f)^2\, \mathcal{S}_2}{\sigma^2} \\
J_{\nu, \nu}
    & \dfrac{8\pi^2 |\tilde{\alpha}|^2 N T_S^2\, Q_2}{\sigma^2}
    & \dfrac{8\pi^2 |\tilde{\alpha}|^2}{\sigma^2}\!\left[T_S^2 Q_2\, \mathcal{S}_0 + 2 T_S Q_1\, \mathrm{Re}\{\mathcal{T}_0\} + M\, \mathcal{U}\right]
    & \dfrac{8\pi^2 |\tilde{\alpha}|^2 N T_S^2\, Q_2}{\sigma^2}
    & \dfrac{8\pi^2 |\tilde{\alpha}|^2}{\sigma^2}\!\left[T_S^2 Q_2\, \mathcal{S}_0 + 2 T_S Q_1\, \mathrm{Re}\{\mathcal{T}_0\} + M\, \mathcal{U}\right] \\
J_{\tau, \nu}
    & -\dfrac{8\pi^2 |\tilde{\alpha}|^2 \Delta f\, T_S\, P_1 Q_1}{\sigma^2}
    & -\dfrac{8\pi^2 |\tilde{\alpha}|^2 \Delta f}{\sigma^2}\!\left[T_S Q_1\, \mathcal{S}_1 + M\, \mathrm{Re}\{\mathcal{T}_1\}\right]
    & -\dfrac{8\pi^2 |\tilde{\alpha}|^2 \Delta f\, T_S\, P_1 Q_1}{\sigma^2}
    & -\dfrac{8\pi^2 |\tilde{\alpha}|^2 \Delta f}{\sigma^2}\!\left[T_S Q_1\, \mathcal{S}_1 + M\, \mathrm{Re}\{\mathcal{T}_1\}\right] \\
J_{\tilde{\alpha}_R, \tau}
    & \dfrac{4\pi\, \tilde{\alpha}_I\, M \Delta f\, P_1}{\sigma^2}
    & \dfrac{4\pi\, \tilde{\alpha}_I\, M \Delta f\, \mathcal{S}_1}{\sigma^2}
    & \dfrac{4\pi\, \tilde{\alpha}_I\, M \Delta f\, P_1}{\sigma^2}
    & \dfrac{4\pi\, \tilde{\alpha}_I\, M \Delta f\, \mathcal{S}_1}{\sigma^2} \\
J_{\tilde{\alpha}_R, \nu}
    & -\dfrac{4\pi\, \tilde{\alpha}_I\, N T_S\, Q_1}{\sigma^2}
    & -\dfrac{4\pi}{\sigma^2}\, \mathrm{Im}\!\left\{\tilde{\alpha}\!\left[T_S Q_1\, \mathcal{S}_0 + M\, \mathcal{T}_0\right]\right\}
    & -\dfrac{4\pi\, \tilde{\alpha}_I\, N T_S\, Q_1}{\sigma^2}
    & -\dfrac{4\pi}{\sigma^2}\, \mathrm{Im}\!\left\{\tilde{\alpha}\!\left[T_S Q_1\, \mathcal{S}_0 + M\, \mathcal{T}_0\right]\right\} \\
J_{\tilde{\alpha}_I, \tau}
    & -\dfrac{4\pi\, \tilde{\alpha}_R\, M \Delta f\, P_1}{\sigma^2}
    & -\dfrac{4\pi\, \tilde{\alpha}_R\, M \Delta f\, \mathcal{S}_1}{\sigma^2}
    & -\dfrac{4\pi\, \tilde{\alpha}_R\, M \Delta f\, P_1}{\sigma^2}
    & -\dfrac{4\pi\, \tilde{\alpha}_R\, M \Delta f\, \mathcal{S}_1}{\sigma^2} \\
J_{\tilde{\alpha}_I, \nu}
    & \dfrac{4\pi\, \tilde{\alpha}_R\, N T_S\, Q_1}{\sigma^2}
    & \dfrac{4\pi}{\sigma^2}\, \mathrm{Re}\!\left\{\tilde{\alpha}\!\left[T_S Q_1\, \mathcal{S}_0 + M\, \mathcal{T}_0\right]\right\}
    & \dfrac{4\pi\, \tilde{\alpha}_R\, N T_S\, Q_1}{\sigma^2}
    & \dfrac{4\pi}{\sigma^2}\, \mathrm{Re}\!\left\{\tilde{\alpha}\!\left[T_S Q_1\, \mathcal{S}_0 + M\, \mathcal{T}_0\right]\right\} \\
\hline
\end{array}
$}
\caption{Expected single-target Fisher information matrix entries
under i.i.d.\ unit-energy QPSK signalling. Columns 1 and 3 (no-ICI
cases) are deterministic in the symbols; columns 2 and 4 (ICI cases)
are expected over the random data. The four off-diagonal entries with
no \(\nu\)-derivative differ between the no-ICI and ICI columns only
through the substitution \(P_a\to\mathcal{S}_a\); the four entries
that carry a \(\nu\)-derivative additionally pick up the
\(\mathcal{T}_a,\mathcal{U}\) corrections inherited from the kernel
\(I\)'s own \(\nu\)-dependence.}
\label{tab:fim_compare}
\end{table}

\section{Equivalence between OFDM and OTFS under QPSK}

The most striking observation in Table~\ref{tab:fim_compare} is that
columns 1 and 3 coincide entry by entry, as do columns 2 and 4. That
is, under i.i.d.\ unit-energy QPSK signalling
\begin{equation}
    \boxed{\;
    \mathrm{E}\{\bm{J}^{\text{OFDM, no ICI}}\} =
    \mathrm{E}\{\bm{J}^{\text{OTFS, no OFDM-ICI}}\},
    \qquad
    \mathrm{E}\{\bm{J}^{\text{OFDM, ICI}}\} =
    \mathrm{E}\{\bm{J}^{\text{OTFS, OFDM-ICI}}\}.
    \;}
\end{equation}
The reason is structural: the ISFFT and SFFT are unitary linear
transforms relating \(X[n,m]\) and \(x[l,k]\) one-to-one. For data
that are jointly circularly symmetric and i.i.d.\ unit-energy --- a
property the unit-energy QPSK ensemble enjoys --- the second-order
statistics that determine the FIM are invariant under this unitary
change of basis. Equivalently, the OTFS observation
\(y[l',k']\) is the SFFT of the OFDM observation \(Y[n,m]\); applying
a unitary transform to a Gaussian observation preserves its
information about \(\bm{\eta}\). The two grid geometries (time--frequency
vs.\ delay--Doppler) are therefore informationally equivalent under
random QPSK loading and the FIM is identical entry by entry.

Two corollaries follow:
\begin{itemize}
    \item \emph{CRBs match.} Any Cram\'er--Rao bound computed from the
          OFDM no-ICI FIM is also the CRB for the no-OFDM-ICI OTFS
          model; likewise the OFDM ICI CRB equals the OFDM-ICI OTFS
          CRB.
    \item \emph{Sensing performance is data-dependent, not
          waveform-dependent.} The choice between placing pilots in
          \((n,m)\) or \((l,k)\) does not change estimator variance
          when data are i.i.d.\ unit-energy QPSK; what matters is the
          ICI-vs.-no-ICI receiver model and the average squared
          symbol energy.
\end{itemize}
The equivalence breaks down if the data depart from i.i.d.\
unit-energy (e.g.\ pilot patterns concentrating energy on specific
grid points, or signalling alphabets with non-constant modulus). In
those cases the OFDM and OTFS FIMs can differ even within the same
ICI model.

\section{No-ICI vs.\ ICI under QPSK}

Since the OFDM/OTFS pairs agree, the remaining comparison reduces to
two distinct expected FIMs --- one without OFDM-side ICI and one with.
Their entry-by-entry differences are:

\paragraph{Entries with no \(\nu\)-derivative
(\(J_{\tilde{\alpha}_R\tilde{\alpha}_R}\),
\(J_{\tilde{\alpha}_I\tilde{\alpha}_I}\),
\(J_{\tau\tau}\),
\(J_{\tilde{\alpha}_R\tau}\),
\(J_{\tilde{\alpha}_I\tau}\)).}
These differ only through the single substitution
\(P_a \to \mathcal{S}_a\). The kernel \(I\) does not appear separately;
it merely smears the discrete index moments \(P_a = \sum_{n} n^a\) into
the lag-summed quadratic moments
\(\mathcal{S}_a = \sum_{n,n'} (n')^a |I(n'-n)|^2\). When \(|I|^2\) is
sharply concentrated at the origin (low normalized Doppler \(\nu T\)),
\(\mathcal{S}_a \approx P_a \cdot \sum_q |I(q)|^2\) and the entries
scale similarly to their no-ICI counterparts up to a real positive
constant.

\paragraph{Entries with a \(\nu\)-derivative
(\(J_{\nu\nu}\),
\(J_{\tau\nu}\),
\(J_{\tilde{\alpha}_R\nu}\),
\(J_{\tilde{\alpha}_I\nu}\)).}
These pick up extra terms in \(\mathcal{T}_a\) and \(\mathcal{U}\)
that have no analog in the no-ICI form. The mechanism is that \(I\)
itself depends on \(\nu\); differentiating with respect to \(\nu\)
therefore excites two contributions:
\begin{enumerate}
    \item the symbol-time phase rotation \(e^{j2\pi mT_S\nu}\), which
          carries through to the no-ICI expressions via the
          \(Q_a = \sum_m m^a\) moments, and
    \item the per-subcarrier kernel response \(\partial I/\partial\nu
          = j2\pi I'\), which has no counterpart when \(I\) is treated
          as a fixed delta.
\end{enumerate}
Because \(\mathcal{T}_0\) and \(\mathcal{T}_1\) are generically
complex, the gain--Doppler cross-entries
(\(J_{\tilde{\alpha}_R\nu}\), \(J_{\tilde{\alpha}_I\nu}\)) mix
\(\tilde{\alpha}_R\) and \(\tilde{\alpha}_I\) through the
\(\mathrm{Im}\{\tilde{\alpha}\,(\cdot)\}\) and
\(\mathrm{Re}\{\tilde{\alpha}\,(\cdot)\}\) wrappers, whereas in the
no-ICI case each entry isolates a single quadrature of
\(\tilde{\alpha}\) cleanly.

\paragraph{Parameter dependence.} Neither FIM depends on the true
delay \(\tau\) --- its phase factors cancel pairwise in both rows of
the table. The Doppler dependence differs sharply:
\begin{itemize}
    \item \emph{No ICI (columns 1, 3)}: the FIM is \emph{independent} of
          the true \(\nu\). All \(\nu\)-phases \(e^{j2\pi mT_S\nu}\)
          drop out under \(|\cdot|^2\) operations, and \(I\) is fixed
          to \(\delta\). The CRB is therefore operating-point free
          --- the same numeric value applies whether the target sits
          on the grid or moves at any Doppler.
    \item \emph{With ICI (columns 2, 4)}: the FIM depends on the
          true \(\nu\) through every occurrence of
          \(\mathcal{S}_a,\mathcal{T}_a,\mathcal{U}\), since these are
          quadratic forms of \(I\) and \(I'\) which both depend on
          \(\nu\). The CRB is therefore operating-point dependent and
          generally degrades as \(\nu T\) grows (broader \(I\), more
          off-diagonal leakage).
\end{itemize}

\paragraph{Asymptotic scaling.}
For large \((N, M)\): \(P_a, Q_a, \mathcal{S}_a \sim N^{a+1}, M^{a+1},
N^{a+1}\) respectively, while \(\mathcal{T}_a \sim N^{a+1}\) and
\(\mathcal{U}\sim N\). Substituting,
\begin{equation}
    J_{\tilde{\alpha}\tilde{\alpha}} \sim NM,
    \quad
    J_{\tau\tau} \sim N^3 M,
    \quad
    J_{\nu\nu} \sim N M^3
\end{equation}
in all four cases --- the dominant grid scaling is unchanged by the
choice of waveform (OFDM vs OTFS) or by the inclusion of ICI. The ICI
signature lives in the constants of proportionality (set by the shape
of \(I\)) and in the sub-leading \(\mathcal{T}_a,\mathcal{U}\)
corrections that become non-negligible at modest \(M\) or when the
operating \(\nu T\) is large enough to spread \(I\) noticeably across
subcarriers.

\paragraph{Limit \(I\to\delta\).} Setting
\(I(n)=\delta[n]\) and \(I'(n)=0\) gives \(\mathcal{S}_a\to P_a\),
\(\mathcal{T}_a\to 0\), \(\mathcal{U}\to 0\), so columns 2 and 4
collapse onto columns 1 and 3 respectively. The four-column table
then degenerates into a single column, confirming internal
consistency.

\section{Multi-target case}

For \(P>1\) targets, all four FIMs share the block structure
\begin{equation}
    \bm{J}(\bm{\eta}) = \begin{bmatrix}
        \bm{J}_{1,1} & \cdots & \bm{J}_{1,P} \\
        \vdots & \ddots & \vdots \\
        \bm{J}_{P,1} & \cdots & \bm{J}_{P,P}
    \end{bmatrix},\qquad
    \bm{J}_{q,p} = \bm{J}_{p,q}^T,\qquad
    \bm{\eta}_p = [\tilde{\alpha}_{p,R},\tilde{\alpha}_{p,I},\tau_p,\nu_p]^T,
\end{equation}
with diagonal blocks \(\bm{J}_{p,p}\in\mathbb{R}^{4\times 4}\) carrying
the per-target information and off-diagonal blocks \(\bm{J}_{p,q}\)
\((p\neq q)\) carrying cross-target coupling.

\subsection*{Diagonal blocks}

The diagonal block \(\bm{J}_{p,p}\) equals the single-target FIM of
the previous section with all parameters indexed by \(p\) and with
\(\tilde{\alpha}\to\tilde{\alpha}_p\). All four cases therefore inherit
the diagonal-block landscape of Section~2 verbatim:
\begin{equation}
    \mathrm{E}\{\bm{J}_{p,p}^{\text{OFDM, no-ICI}}\}
    = \mathrm{E}\{\bm{J}_{p,p}^{\text{OTFS, no-OFDM-ICI}}\},\qquad
    \mathrm{E}\{\bm{J}_{p,p}^{\text{OFDM, ICI}}\}
    = \mathrm{E}\{\bm{J}_{p,p}^{\text{OTFS, OFDM-ICI}}\},
\end{equation}
and Table~\ref{tab:fim_compare} applies to each diagonal block under
the substitutions \(\tilde{\alpha}\to\tilde{\alpha}_p\),
\(\tau\to\tau_p\), \(\nu\to\nu_p\), \(I\to I_p\). The four-way
comparison again collapses to two distinct entries (no-ICI vs.\ ICI).

\subsection*{Cross-target moments}

For the off-diagonal blocks define the inter-target separations
\(\Delta\tau_{pq} := \tau_q - \tau_p\),
\(\Delta\nu_{pq} := \nu_q - \nu_p\), the Doppler-side Dirichlet sum
\begin{equation}
    F_b(z,K) := \sum_{k=0}^{K-1} k^b\, e^{j2\pi kz},
\end{equation}
and the cross-target kernel-summed moments
\begin{align}
    \mathcal{S}_a^{p,q} &:= \sum_{n,n'=0}^{N-1} (n')^a\,
        I_p^*(n'-n)\, I_q(n'-n)\,
        e^{-j2\pi n'\,\Delta f\,\Delta\tau_{pq}},\quad a=0,1,2, \\
    \mathcal{T}_a^{p,q} &:= \sum_{n,n'=0}^{N-1} (n')^a\,
        I_p^*(n'-n)\, I_q'(n'-n)\,
        e^{-j2\pi n'\,\Delta f\,\Delta\tau_{pq}},\quad a=0,1, \\
    \check{\mathcal{T}}_a^{p,q} &:= \sum_{n,n'=0}^{N-1} (n')^a\,
        (I_p')^*(n'-n)\, I_q(n'-n)\,
        e^{-j2\pi n'\,\Delta f\,\Delta\tau_{pq}}, \\
    \mathcal{U}^{p,q} &:= \sum_{n,n'=0}^{N-1}
        (I_p')^*(n'-n)\, I_q'(n'-n)\,
        e^{-j2\pi n'\,\Delta f\,\Delta\tau_{pq}},
\end{align}
which generalise \(\mathcal{S}_a,\mathcal{T}_a,\mathcal{U}\) of
Section~1. At \(p=q\) the phase factor disappears and
\(\mathcal{S}_a^{p,p}=\mathcal{S}_a\),
\(\mathcal{T}_a^{p,p}=\mathcal{T}_a\),
\(\check{\mathcal{T}}_a^{p,p}=\mathcal{T}_a^*\),
\(\mathcal{U}^{p,p}=\mathcal{U}\). In the limit \(I_p\to\delta\),
\(\mathcal{S}_a^{p,q}\to F_a(-\Delta f\,\Delta\tau_{pq},N)\) and
\(\mathcal{T}_a^{p,q},\check{\mathcal{T}}_a^{p,q},\mathcal{U}^{p,q}\to 0\),
so the ICI cross-target moments reduce to the no-ICI delay-side
ambiguity function. We abbreviate
\(F_a^{(\tau)} := F_a(-\Delta f\,\Delta\tau_{pq},N)\) and
\(F_b^{(\nu)} := F_b(T_S\,\Delta\nu_{pq},M)\) throughout.

\subsection*{Off-diagonal block comparison}

Under QPSK signalling, the (I)SFFT-unitarity argument that gave
\(\mathrm{E}\{\bm{J}^{\text{OFDM}}\}=\mathrm{E}\{\bm{J}^{\text{OTFS}}\}\)
for each diagonal block applies equally to the off-diagonal blocks.
Concretely:
\begin{equation}
    \mathrm{E}\{\bm{J}_{p,q}^{\text{OFDM, no ICI}}\}
    = \mathrm{E}\{\bm{J}_{p,q}^{\text{OTFS, no OFDM-ICI}}\},
    \qquad
    \mathrm{E}\{\bm{J}_{p,q}^{\text{OFDM, ICI}}\}
    = \mathrm{E}\{\bm{J}_{p,q}^{\text{OTFS, OFDM-ICI}}\}
    \quad\text{for every } p,q.
\end{equation}
The four-case off-diagonal comparison therefore collapses to a
two-column table (no-ICI form vs.\ ICI form), shown in
Table~\ref{tab:fim_cross}. The six rows cover one entry from each
``family'' (gain--gain, gain--delay, gain--Doppler, delay--delay,
Doppler--Doppler, delay--Doppler).

\begin{table}[h]
\centering
\small
\renewcommand{\arraystretch}{1.8}
\[
\begin{array}{l|c|c}
\hline
\text{Entry } J_{p,q,(a,b)}
    & \text{No OFDM-ICI}
    & \text{With OFDM-ICI} \\
\hline
J_{p,q,\tilde{\alpha}_R\tilde{\alpha}_R}
    & \dfrac{2}{\sigma^2}\,\mathrm{Re}\{F_0^{(\tau)} F_0^{(\nu)}\}
    & \dfrac{2}{\sigma^2}\,\mathrm{Re}\{\mathcal{S}_0^{p,q}\, F_0^{(\nu)}\} \\
J_{p,q,\tilde{\alpha}_R\tau}
    & \dfrac{4\pi\,\Delta f}{\sigma^2}\,\mathrm{Im}\{\tilde{\alpha}_q\, F_1^{(\tau)} F_0^{(\nu)}\}
    & \dfrac{4\pi\,\Delta f}{\sigma^2}\,\mathrm{Im}\{\tilde{\alpha}_q\,\mathcal{S}_1^{p,q}\, F_0^{(\nu)}\} \\
J_{p,q,\tilde{\alpha}_R\nu}
    & -\dfrac{4\pi\, T_S}{\sigma^2}\,\mathrm{Im}\{\tilde{\alpha}_q\, F_0^{(\tau)} F_1^{(\nu)}\}
    & -\dfrac{4\pi}{\sigma^2}\,\mathrm{Im}\{\tilde{\alpha}_q(T_S\mathcal{S}_0^{p,q} F_1^{(\nu)} + \mathcal{T}_0^{p,q} F_0^{(\nu)})\} \\
J_{p,q,\tau\tau}
    & \dfrac{8\pi^2 (\Delta f)^2}{\sigma^2}\,\mathrm{Re}\{\tilde{\alpha}_p^*\tilde{\alpha}_q F_2^{(\tau)} F_0^{(\nu)}\}
    & \dfrac{8\pi^2 (\Delta f)^2}{\sigma^2}\,\mathrm{Re}\{\tilde{\alpha}_p^*\tilde{\alpha}_q \mathcal{S}_2^{p,q} F_0^{(\nu)}\} \\
J_{p,q,\tau\nu}
    & -\dfrac{8\pi^2 \Delta f\, T_S}{\sigma^2}\,\mathrm{Re}\{\tilde{\alpha}_p^*\tilde{\alpha}_q F_1^{(\tau)} F_1^{(\nu)}\}
    & -\dfrac{8\pi^2 \Delta f}{\sigma^2}\,\mathrm{Re}\{\tilde{\alpha}_p^*\tilde{\alpha}_q (T_S\mathcal{S}_1^{p,q} F_1^{(\nu)} + \mathcal{T}_1^{p,q} F_0^{(\nu)})\} \\
\hline
\end{array}
\]
\caption{Five representative entries of the expected off-diagonal
block \(\bm{J}_{p,q}\) (\(p\neq q\)) under i.i.d.\ unit-energy QPSK
signalling. The OFDM and OTFS columns coincide for each ICI model
(four-case comparison collapses to two), so only the two distinct
forms are listed. Notation:
\(F_a^{(\tau)} := F_a(-\Delta f\,\Delta\tau_{pq},N)\) and
\(F_b^{(\nu)} := F_b(T_S\,\Delta\nu_{pq},M)\). The
\(\nu\nu\)-entry is displayed separately in
equation~\eqref{eq:Jpq_nunu} due to its length; the remaining 10
off-diagonal entries follow from the same template by swapping
\(\mathrm{Re}\leftrightarrow\mathrm{Im}\), conjugating
\(\tilde{\alpha}\) factors, and replacing
\(\mathcal{T}\leftrightarrow\check{\mathcal{T}}\) where the role of
\(p\) and \(q\) is reversed.}
\label{tab:fim_cross}
\end{table}

The \(\nu\nu\)-entry, which carries all three of the ICI-induced
moments, is too wide for the table; the no-ICI and ICI forms are
\begin{align}
    {[\bm{J}_{p,q}]_{\nu,\nu}^{\text{no ICI}}} &=
        \tfrac{8\pi^2 T_S^2}{\sigma^2}\,
        \mathrm{Re}\!\left\{\tilde{\alpha}_p^*\tilde{\alpha}_q\,
            F_0^{(\tau)}\, F_2^{(\nu)}\right\}, \\
    {[\bm{J}_{p,q}]_{\nu,\nu}^{\text{ICI}}} &=
        \tfrac{8\pi^2}{\sigma^2}\,
        \mathrm{Re}\!\left\{\tilde{\alpha}_p^*\tilde{\alpha}_q\!
            \left(T_S^2\,\mathcal{S}_0^{p,q}\, F_2^{(\nu)}
                + T_S\bigl(\mathcal{T}_0^{p,q} + \check{\mathcal{T}}_0^{p,q}\bigr)\, F_1^{(\nu)}
                + \mathcal{U}^{p,q}\, F_0^{(\nu)}\right)\right\}.
\label{eq:Jpq_nunu}
\end{align}

\subsection*{Multi-target observations}

\paragraph{Equivalence persists block-by-block.} The same
(I)SFFT-unitarity argument that proved
\(\mathrm{E}\{\bm{J}^{\text{OFDM}}\} =
\mathrm{E}\{\bm{J}^{\text{OTFS}}\}\) for single targets applies
unchanged for each block:
\begin{equation}
    \boxed{\;
    \mathrm{E}\{\bm{J}_{p,q}^{\text{OFDM, no ICI}}\} =
    \mathrm{E}\{\bm{J}_{p,q}^{\text{OTFS, no OFDM-ICI}}\},
    \qquad
    \mathrm{E}\{\bm{J}_{p,q}^{\text{OFDM, ICI}}\} =
    \mathrm{E}\{\bm{J}_{p,q}^{\text{OTFS, OFDM-ICI}}\}
    \quad\forall\, p,q.
    \;}
\end{equation}
The whole \(4P\times 4P\) FIM is therefore identical between OFDM and
OTFS under i.i.d.\ unit-energy QPSK loading; only the ICI/no-ICI
distinction matters.

\paragraph{No-ICI vs.\ ICI in the cross blocks.} As in the
single-target case, the ICI columns differ from the no-ICI columns
through two patterns:
\begin{itemize}
    \item \emph{Entries with no \(\nu\)-derivative} (the
          \(\tilde{\alpha}_R\tilde{\alpha}_R\),
          \(\tilde{\alpha}_R\tau\) and \(\tau\tau\) rows of
          Table~\ref{tab:fim_cross}) inherit a single substitution
          \(F_a^{(\tau)} \to \mathcal{S}_a^{p,q}\). The Doppler-side
          factor \(F_b^{(\nu)}\) is unchanged. Crucially,
          \(\mathcal{S}_a^{p,q}\) is \emph{not} a simple Dirichlet
          sum: it is a weighted lag correlation of \(I_p\) and
          \(I_q\) carrying the inter-target delay phase. Coincident
          targets (\(\tau_p = \tau_q\)) give the real
          single-target \(\mathcal{S}_a\); well-separated targets
          give a complex value that decays with \(\Delta\tau_{pq}\)
          but at a rate set by the \emph{full kernel width} of
          \(|I|\), not the bin width \(1/(N\Delta f)\).
    \item \emph{Entries with a \(\nu\)-derivative} (the
          \(\tilde{\alpha}_R\nu\) and \(\tau\nu\) rows of
          Table~\ref{tab:fim_cross}, plus the \(\nu\nu\)-entry
          \eqref{eq:Jpq_nunu}) additionally pick up the cross-target
          \(\mathcal{T}_a^{p,q}, \check{\mathcal{T}}_a^{p,q},
          \mathcal{U}^{p,q}\) moments, which have no analog in the
          no-ICI form. These produce extra cross-target coupling
          terms whenever \(I\) has appreciable Doppler bandwidth, and
          they survive even when \(\Delta\nu_{pq}=0\) (i.e.\ two
          targets at the same Doppler but different delays still
          couple through \(\mathcal{T},\mathcal{U}\) in the ICI
          model).
\end{itemize}

\paragraph{Resolution and block-diagonality.} For all four cases the
off-diagonal block scales with the underlying ambiguity functions:
\begin{itemize}
    \item No-ICI: amplitude controlled by
          \(|F_a^{(\tau)}|\,|F_b^{(\nu)}|\) which decays once
          \(|\Delta\tau_{pq}|\gtrsim 1/(N\Delta f)\) or
          \(|\Delta\nu_{pq}|\gtrsim 1/(MT_S)\) --- the delay/Doppler
          bin widths.
    \item ICI: \(|\mathcal{S}_a^{p,q}|\) decays at the slower
          \emph{kernel-width} rate \(\sim 1/T \sim \Delta f\)
          (because \(|I|\) is itself
          \(\mathrm{sinc}\)-like with width \(\sim\Delta f\)), while
          \(|F_b^{(\nu)}|\) still decays at \(1/(MT_S)\) on the
          Doppler side.
\end{itemize}
In either case, well-separated targets give
\(\bm{J}\approx\mathrm{blkdiag}(\bm{J}_{p,p})\) and the multi-target
CRB factorises into independent per-target CRBs. At sub-resolution
spacings the cross blocks become comparable to the diagonal blocks
and the multi-target CRB inflates as expected.

\paragraph{Limit \(I\to\delta\).} Setting
\(I_p(n)=\delta[n]\) and \(I_p'(n)=0\) for every \(p\) gives
\(\mathcal{S}_a^{p,q}\to F_a^{(\tau)}\),
\(\mathcal{T}_a^{p,q},\check{\mathcal{T}}_a^{p,q},\mathcal{U}^{p,q}\to 0\),
collapsing the ICI column of Table~\ref{tab:fim_cross} entry-by-entry
onto the no-ICI column --- the multi-target analog of the single-target
limit at the end of Section~3.

\section{Summary}

Under the i.i.d.\ unit-energy QPSK assumption, the four single-target
FIMs and the four multi-target FIMs each collapse to \emph{two}
distinct cases:
\begin{itemize}
    \item \textbf{No OFDM-side ICI} (OFDM no-ICI \(=\) OTFS no-OFDM-ICI):
          single-target diagonal blocks are closed-form,
          \(\tau\)- and \(\nu\)-independent, driven by the elementary
          power sums \(P_a, Q_b\) of the grid. Multi-target
          off-diagonal blocks are driven by the delay/Doppler
          ambiguity functions
          \(F_a(-\Delta f\,\Delta\tau_{pq},N)\,F_b(T_S\,\Delta\nu_{pq},M)\),
          which decay at the grid bin widths
          \(1/(N\Delta f)\) and \(1/(MT_S)\) respectively. CRBs are
          deterministic and operating-point-free.
    \item \textbf{With OFDM-side ICI} (OFDM ICI \(=\) OTFS OFDM-ICI):
          single-target diagonal blocks are driven by the
          kernel-summed moments \(\mathcal{S}_a, \mathcal{T}_a,
          \mathcal{U}\), which depend on the true Doppler through
          \(I\) and \(I'\). Multi-target off-diagonal blocks involve
          their cross-target generalisations
          \(\mathcal{S}_a^{p,q}, \mathcal{T}_a^{p,q},
          \check{\mathcal{T}}_a^{p,q}, \mathcal{U}^{p,q}\) on the
          delay side and \(F_b(T_S\,\Delta\nu_{pq},M)\) on the Doppler
          side. CRBs are operating-point dependent and recover the
          no-ICI form continuously as \(\nu T\to 0\).
\end{itemize}
The OFDM/OTFS labels distinguish only the data-domain interpretation
(time--frequency vs.\ delay--Doppler); the information content is the
same for both single-target and multi-target estimation so long as the
symbols are i.i.d.\ unit-energy. To break this equivalence one must
depart from i.i.d.\ unit-energy symbols --- pilots, non-constant-modulus
alphabets, or correlated training sequences --- in which case the OFDM
and OTFS columns can separate and the choice of waveform begins to
matter for parameter estimation.

\end{document}
